\ssscp{Vowel in final syllable}{14-04-03-01}
Beginning with the final vowel,
most of the data in languages in Wallacea (in Tables \ref{tab:14-07a}, \ref{tab:14-07b}, and \ref{tab:14-07c}
and appendix \ref{ch:SurMar}) can be accounted for by *ə.
\tsc{Sula-Buru} has regular *ə > \it{e} (\trf{09-01}),
\tsc{Seram-Tanimbar-Bomberai} has regular *ə > \it{a} in final
syllables (\srf{11-01-02}), and \ili{Banda} has *ə > \it{e} after penultimate \it{e} (\srf{12-05}).

Much of \tsc{\ili{Ambon}-Seram} also has final *ə > \it{e}.
(see: \trf{09-01}, *Raməs ‘press’ [\trf{10-02}], *kədəŋ ‘stand’,
*daləm ‘inside’, *waRəj ‘vine,
creeper’, **niwəR ‘coconut’ [all \trf{10-04}],
and *quləj ‘worm’ [\trf{10-07}]).
\ili{Nusalaut} and \ili{Saparua} have final *ə > \it{o} meaning that we have
to posit irregular *ə > \it{e} for \it{makel} ‘cuscus’ for these languages.
Two other \tsc{\ili{Ambon}-Seram} words show irregularities in the
final vowel: South \ili{Nuaulu} \it{mar\tbr{a}n-e} ‘cuscus’ has *ə
> \it{a} against usual final-syllable *ə > \it{e} (example \qf{10-57},
\srf{10-03-03}), and \ili{Liana-Seti} \it{ma\tbr{i}r-a} has *ə > \it{i}
against usual final-syllable *ə > \it{e, a}.

\citet[264]{Schapper2011} reconstructs **medaR ‘cuscus’ for Proto-\ili{Aru}.
This has final-syllable *ə > **a,
which is one of the outcomes for final-syllable *ə reconstructed
by \citet{Nivens2017}, as the data in \trf{14-08} show.
(Final-syllable *ə > **i
and *ə > **u are also attested in \tsc{Aru}.) The East \ili{Kola} term
given by \citet{Nivens2017} is \it{moduh} with final /u/ which
cannot be regularly accounted for by final *a or *ə.

\begin{table}[h]
%\begin{adjustbox}{width=\textwidth}
	\centering\tcp{Proto-\ili{Aru} *ə > **a}{14-08}
	\st{2pt}\begin{tabular}{l|lll@{}lll}\lsptoprule
local gl.	&	{\hr}hear	&	{\hr}sever	&	{\hr}wide	&	{\hr}caterpillar	&	{\hr}swallow	&	{\hr}black\\
PMP	&	\hr*dəŋ\tbr{ə}R	&	\hr*təkt\tbr{ə}k	&	\hr*lab\tbr{ə}R	&	\hr*qul\tbr{ə}j	&	\hr*təl\tbr{ə}n	&	\hr*ma-qit\tbr{ə}m\\	\midrule
Proto-\ili{Aru}	&	**-reŋ\tbr{a}R	&	**-tet\tbr{a}k	&	**laɸ\tbr{a}r-	&	**(wy)il\tbr{a}R	&	**-tel\tbr{a}(m)	&	\\
\ili{Ujir}	&	{\hr\hr}-reŋ\tbr{a}y	&	{\hr\hr}-tet\tbr{a}k	&	{\hr\hr}laɸ\tbr{i}	&	{\hr\hr}wil{\tl}wil\tbr{a}y	&		\mc{2}{r}{{\hr\hr}aŋamet\tbr{a}m}\\
East \ili{Kola}	&	{\hr\hr}-ren\tbr{a}	&	{\hr\hr}-tet\tbr{a}	&	{\hr\hr}lab\tbr{i}h	&	{\hr\hr}yil\tbr{a}h	&	{\hr\hr}-tel\tbr{e}	&	\\
\ili{Dobel}	&	{\hr\hr}-reŋ\tbr{i}n	&	{\hr\hr}te-taʔ\tbr{u}	&		&	{\hr\hr}s{\tl}sil\tbr{a}r	&	{\hr\hr}-yel\tbr{i}m	&	\\
SW \ili{Tarangan}	&	{\hr\hr}-luŋ\tbr{u}r	&	{\hr\hr}-set\tbr{a}k	&	{\hr\hr}laɸ\tbr{a}r-di	&	{\hr\hr}ʤil\tbr{a}r	&	{\hr\hr}-sel\tbr{a}	&	\\
		\lspbottomrule
	\end{tabular}
%\end{adjustbox}
\end{table}

\tsc{Timor-Babar}
and \tsc{\ili{Central} Timor} languages usually have final *ə > \it{a}
thus accounting for the final vowels in nearly all these languages.\footnote{
\ili{Makuva}
(\srf{04-06-01}) and Waima{\Q}a have regular final *ə > \it{e}.}
The final vowel
in \ili{Meto} (\ili{Ketun}, Nai'oni' village) \it{ʔmakuʔ} or Fatule{\Q}u \it{maukuʔ} ‘cuscus’
is not regular from either *ə or *a,
as both regularly become \it{a} in final syllables in \ili{Meto} \citep[55f, 73]{Edwards2021}.
Both these have variant forms with final \it{a:} \ili{Ketun} (Taloetan
village) \it{ʔmaukaʔ,} Fatule{\Q}u \it{maukaʔ} indicating a local
irregular shift of **a \it{> u}.
There is another \ili{Meto} term for cuscus,
**arum \citep[91f]{Edwards2021} and irregular **a > \it{u}
in \it{ʔmakuʔ} (\ili{Ketun}, Nai'oni' village), \it{maukuʔ} (Fatule{\Q}u) may be influenced
by the final vowel in this other term.

\tsc{\ili{Muna}-Buton} languages usually have *ə > \it{o},
including in final syllables (found in 10/15 examples checked),
meaning that final \it{e} in \it{mante} ‘bear cuscus’ may be unexpected.
However, PMP *R/*j became Proto-Nuclear \tsc{\ili{Muna}-Buton} **y
\citep[91]{vandenBerg2003},
which was then often lost but caused raising of an adjacent vowel.
Note in \trf{14-09} below that the reflexes of word-final *əR{\#} are not universally \it{o}.
In fact, *əR > \it{e} /{\gap}{\#} seems to be regular in some languages,
as shown in the shaded area.
Thus, final *əR in *mantəR ‘cuscus’ could yield \it{e},
with the associated loss of final *R.
The data in \trf{14-09} are gleaned from David Mead’s word lists.

\begin{table}[h]
%\begin{adjustbox}{width=\textwidth}
	\centering\tcp{\tsc{\ili{Muna}-Buton} *R > **y and *əR > \it{e} /{\gap}{\#}}{14-09}
	\st{2pt}\begin{tabular}{l|llllll}\lsptoprule
*gloss	&	blood	&	thorn	&	vein	&	hear	&	penis	&	wide\\
PMP	&	*da\tbr{R}aq	&	*du\tbr{R}i	&	*u\tbr{R}at	&	*dəŋ\tbr{əR}	&	*las\tbr{əR}	&	*ma-lab\tbr{əR}\\	\midrule
\ili{Wakole}	&	{\hr}xea	&	{\hr}xui	&	{\hr}ue	&	{\hr}pindoŋo	&		&	{\hr}moleβa\\
\ili{Wolowa}	&	{\hr}xea	&	{\hr}xui	&	{\hr}ue	&	{\hr}pindoŋo	&		&	{\hr}moleβa\\
\ili{Kumbewaha}	&	{\hr}xea	&	{\hr}xui	&	{\hr}uβe-no xea	&	{\hr}pindoŋo	&		&	\\
\ili{Kambowa}	&	{\hr}rea	&	{\hr}rii	&	{\hr}ie	&	{\hr}ro-reɗen\tbr{e}	{\cdr}&	{\hr}leh\tbr{e}	{\cdr}&	{\hr}no-moleβ\tbr{e}{\cdr}\\
\ili{Kioko}	&	{\hr}xea	&	{\hr}xii	&	{\hr}ie	&	{\hr}no-xeɗen\tbr{e}	{\cdr}&	{\hr}leʔ\tbr{e}	{\cdr}&	{\hr}no-moleβ\tbr{e}{\cdr}\\
\ili{Pancana}	&	{\hr}rea	&	{\hr}rii	&	{\hr}iʔei	&	{\hr}reɗen\tbr{e}, reden\tbr{e}	{\cdr}&		{\cdr}&	{\hr}moleβ\tbr{e}{\cdr}\\
\ili{Kaimbulawa}	&	{\hr}eʕa	&	{\hr}kiʕi	&	{\hr}uye	&	{\hr}aɗen\tbr{e}	{\cdr}&		{\cdr}&	{\hr}moleβ\tbr{e}{\cdr}\\
\ili{Busoa}	&	{\hr}ɣea	&	{\hr}ɣii	&	{\hr}uβa	&		&		&	{\hr}maewa\\
\ili{Kaisabu}	&	{\hr}xea	&	{\hr}kaxii	&	{\hr}uβa	&	{\hr}pendoŋo	&		&	{\hr}moleβa\\
\ili{Cia-Cia}	&	{\hr}rea	&	{\hr}rui	&	{\hr}uβa, ua	&	{\hr}pindoŋo, rodoŋo	&	{\hr}laʔa	&	{\hr}moleβa\\
		\lspbottomrule
	\end{tabular}
%\end{adjustbox}
\end{table}
